\documentclass[a4paper,10pt]{article}
\usepackage[left=1.5cm,top=2cm,right=2.5cm,bottom=4cm]{geometry} 
\usepackage[utf8]{inputenc}
\usepackage[activeacute,spanish]{babel}
\usepackage{multirow}	% Para poder unir filas en las tablas
\usepackage{colortbl}	% Para colorear tablas
\usepackage{latexsym,amssymb,oldlfont,amsmath}
\usepackage{graphicx,graphics}
\usepackage{multicol}
\usepackage{color}
\usepackage{listings}
\usepackage{lastpage}
\usepackage{enumitem}
\usepackage{dsfont} %para poner numero reales
\usepackage{tikz}
\usepackage{tkz-euclide}
\usetkzobj{all}
\usetikzlibrary{calc,angles,positioning,intersections,quotes,decorations.markings}



%\usepackage{fontspec}
%\setmainfont{Times New Roman}
\newcommand{\R}{{\mathbb R}}
\newcommand{\si}{{\mathcal S}^n}
\renewcommand{\theequation}{\thesection.\arabic{equation}} % Cambia formato salida numeracion equation
\makeatletter\@addtoreset {equation}{section} \makeatother % Tambien se usa
\renewcommand{\baselinestretch}{1.3}
\usepackage{amssymb}
\usepackage{hyperref}
\definecolor{midnightblue}{rgb}{0.1,0.1,0.44}


\newcommand{\dis}{\displaystyle}
\providecommand{\abs}[1]{\lvert#1\rvert}
\def\Z{\mathbb{Z}}
\def\N{\mathbb{N}}
\def\R{\mathbb{R}}
\def\P{\mathbb{P}}
\def\M{\mathcal{M}}
\def\sin{\operatorname{sen}}
\def\dim{\operatorname{dim}}
\def\tr{\operatorname{tr}}
\def\Im{\operatorname{Im}}
\def\Ker{\operatorname{Ker}}



\usepackage{etoolbox,fancyhdr,xcolor}
\fancyhf{} % sets both header and footer to nothing
\renewcommand{\headrulewidth}{0pt}
%\renewcommand{\footrulewidth}{0.5pt}
\pagestyle{fancy}
\headheight=60pt %para cambiar el tamaño del encabezadox
\fancyhead[L] %la "L" indica a la izquierda
{
\begin{minipage}{6cm}
\includegraphics[width=0.9\textwidth]{logo_udp.eps}
\end{minipage}
}

%\fancyhead[C]
%{
%\textbf{Cálculo II - Guía 1}
%}


%\fancyhead[R] %la "L" indica a la izquierda
%{

%}
%\begin{minipage}{4cm}
%\small %tamaño de letra pequeño
%{
%\textsf
%{
%Universidad de Chile\\
%...}}
%\end{minipage}}




\fancyfoot[C]{\thepage\ de \pageref{LastPage}}
%Página \thepage\ de \pageref{LastPage}}% Páginas Pares numeradas a la derecha
%\fancyfoot[LO]{\thepage} % Páginas Impares numeradas a la izquierda

\title{AyG-Guia5}

\begin{document}

\begin{center}
\emph{Álgebra y Geometría} \\
\emph{Guía Nº5 Ejercicios Resueltos} \\
\emph{Primer Semestre 2017} \\
\emph{Instituto de Ciencias Básicas}
\end{center}

\begin{center}
\textbf{{\Large Cónicas 1}}
\end{center}

\vspace{2mm}


\begin{enumerate}
\item \textbf{Ejercicio:} Dos lados de un triangulo miden 5 y 10 centímetros, y el angulo que se forma entre ellos de 60º. Determine el perímetro del área del triangulo. \\

\textbf{Solución:}

\begin{center}
\begin{tikzpicture}[scale=0.8]

    % draw the background
    \draw [line width=1.5pt, fill=gray!2] (0,0) -- (60:4) -- (4,0) -- cycle;

    \coordinate[label=left:$A$]  (A) at (0,0);
    \coordinate[label=right:$B$] (B) at (4,0);
    \coordinate[label=above:$C$] (C) at (2,3.464);

    \coordinate[label=below:$10$](c) at ($ (A)!.5!(B) $);
    \coordinate[label=left:$5$] (b) at ($ (A)!.5!(C) $);
    \coordinate[label=right:$x$](a) at ($ (B)!.5!(C) $);

    % angle alpha
    \draw[fill=green!30] (0,0) -- (0:0.75cm) arc (0:60:.75cm);
    \draw (0.35cm,0.25cm) node {$60$};

    % angle beta
    \begin{scope}[shift={(4cm,0cm)}]
        \draw[fill=green!30] (0,0) -- (-180:0.75cm) arc (180:120:0.75cm);
        \draw (150:0.5cm) node {$\beta$};
    \end{scope}

    % angle gamma
    \begin{scope}[shift={(60:4)}]
        \draw[fill=green!30] (0,0) -- (-120:.75cm) arc (-120:-60:.75cm);
        \draw (-90:0.5cm) node {$\gamma$};
    \end{scope}

    % the triangle
    \draw [line width=1.5pt] (A) -- (B) -- (C) -- cycle;
  \end{tikzpicture}
\end{center}


Utilizando el teorema del coseno tenemos que:

\begin{align*}
\dis  a^2  &= b^2 + c^2 - 2*b*c*\cos (\alpha)\\
 \dis  x^2  &= 5^2 + 10^2 - 2*5*10\cos (60)\\
 \dis  x^2  &= 25 + 100 - 100* \frac{1}{2}\\
 \dis  x^2  &= 25 + 100 - 50\\
 \dis  x^2  &= 75\\
 \dis  x  &= \sqrt{75}\\
\end{align*}

\begin{center}
\begin{tikzpicture}[scale=0.8]

    % draw the background
    \draw [line width=1.5pt, fill=gray!2] (0,0) -- (60:4) -- (4,0) -- cycle;

    \coordinate[label=left:$A$]  (A) at (0,0);
    \coordinate[label=right:$B$] (B) at (4,0);
    \coordinate[label=above:$C$] (C) at (2,3.464);

    \coordinate[label=below:$10$](c) at ($ (A)!.5!(B) $);
    \coordinate[label=left:$5$] (b) at ($ (A)!.5!(C) $);
    \coordinate[label=right:$\sqrt{75}$](a) at ($ (B)!.5!(C) $);

    % angle alpha
    \draw[fill=green!30] (0,0) -- (0:0.75cm) arc (0:60:.75cm);
    \draw (0.35cm,0.25cm) node {$60$};

    % angle beta
    \begin{scope}[shift={(4cm,0cm)}]
        \draw[fill=green!30] (0,0) -- (-180:0.75cm) arc (180:120:0.75cm);
        \draw (150:0.5cm) node {$\beta$};
    \end{scope}

    % angle gamma
    \begin{scope}[shift={(60:4)}]
        \draw[fill=green!30] (0,0) -- (-120:.75cm) arc (-120:-60:.75cm);
        \draw (-90:0.5cm) node {$\gamma$};
    \end{scope}

    % the triangle
    \draw [line width=1.5pt] (A) -- (B) -- (C) -- cycle;
  \end{tikzpicture}
\end{center}

Ahora que conocemos los valores de todos los lados del triangulo utilizamos el Teorema de Seno para obtener los ángulos interiores que desconocemos.
\begin{center}
$\dis \frac{\sin(\alpha)}{a} = \frac{\sin(\beta)}{b} = \frac{\sin(\gamma)}{c}$
\end{center}

Buscamos $\beta$:
\begin{align*}
 \dis  \frac{\sin(60)}{\sqrt{75}}  &= \frac{\sin(\beta)}{5}\\
 \dis  5*\sin (60)  &= \sqrt{75}* \sin(\beta)\\
 \dis  \sin (\beta)  &= \frac{5* \frac{\sqrt{3}}{2}}{5\sqrt{3}}\\
 \dis  \sin (\beta)  &= \frac{1}{2}\\
 \dis  \beta  &= 30\\
\end{align*}

Buscamos $\gamma$:
\begin{align*}
 \dis  \frac{\sin(30)}{5}  &= \frac{\sin(\gamma)}{10}\\
 \dis  \sin (\gamma)  &= \frac{10* \frac{1}{2}}{5}\\
 \dis  \sin (\gamma)  &= 1\\
 \dis  \gamma  &= 90\\
\end{align*}

Finalmente obteniendo lo siguiente:

\begin{center}
\begin{tikzpicture}[scale=0.8]

    % draw the background
    \draw [line width=1.5pt, fill=gray!2] (0,0) -- (60:4) -- (4,0) -- cycle;

    \coordinate[label=left:$A$]  (A) at (0,0);
    \coordinate[label=right:$B$] (B) at (4,0);
    \coordinate[label=above:$C$] (C) at (2,3.464);

    \coordinate[label=below:$10$](c) at ($ (A)!.5!(B) $);
    \coordinate[label=left:$5$] (b) at ($ (A)!.5!(C) $);
    \coordinate[label=right:$\sqrt{75}$](a) at ($ (B)!.5!(C) $);

    % angle alpha
    \draw[fill=green!30] (0,0) -- (0:0.75cm) arc (0:60:.75cm);
    \draw (0.35cm,0.25cm) node {$60$};

    % angle beta
    \begin{scope}[shift={(4cm,0cm)}]
        \draw[fill=green!30] (0,0) -- (-180:0.75cm) arc (180:120:0.75cm);
        \draw (150:0.5cm) node {$30$};
    \end{scope}

    % angle gamma
    \begin{scope}[shift={(60:4)}]
        \draw[fill=green!30] (0,0) -- (-120:.75cm) arc (-120:-60:.75cm);
        \draw (-90:0.5cm) node {$90$};
    \end{scope}

    % the triangle
    \draw [line width=1.5pt] (A) -- (B) -- (C) -- cycle;
  \end{tikzpicture}
\end{center}
\newpage

\item \textbf{Ejercicio:} Dos lados de un paralelógramo están dados por las ecuaciones $l_1:x+2y+4=0$ y $l_2:5x+3y-1=0$.\\
Si el punto de intersección de sus diagonales es $E(3,0)$ encuentre las ecuaciones de los otros lados del paralelógramo.\\

\textbf{Solución:}

\begin{center}
\begin{tikzpicture}

%Parallelogram ABCD and its diagonals are drawn.
\coordinate (A) at (0,0);
\coordinate (B) at (105:3);
\coordinate (C) at ($(B) +(5,0)$);
\coordinate (D) at ($(C) +(-75:3)$);

\draw (A) -- (B) -- (C) -- (D) -- cycle;

\draw[name path=path_AC] (A) -- (C);
\draw[name path=path_BD] (B) -- (D);

%The four vertices are labeled.
\node at ($(A)! -2.5mm! (C)$){$A$};
\node at ($(B)! -2.5mm! (D)$){$B$};
\node at ($(C)! -2.5mm! (A)$){$C$};
\node at ($(D)! -2.5mm! (B)$){$D$};

%The intersection of the diagonals of the parallelogram is P.  To place
%the label for P, the line segments are extended 4mm by an invisible
%path command "above" P. The endpoints of these extensions are called
%"label_P_left" and "label_P_right" and the midpoint of the line segment
%between them is called "label_P".  A node command typesets "P" at the
%point 2.5mm from P on the invisible line segment between P and label_P.
\coordinate[name intersections={of=path_AC and path_BD, by={P}}];

\coordinate (label_P_left) at ($(P)!-4mm!(D)$);
\coordinate (label_P_right) at ($(P)!-4mm!(A)$);
\coordinate (label_P) at ($(label_P_left)!0.5!(label_P_right)$);
\node[blue] at ($(P)!3.2mm!(label_P)$){$E(3,0)$};


%Opposite sides of the parallelogram are marked to indicate that they are
%congruent.

\coordinate (mark_for_AD) at ($(A)!0.5!(D)$);
\draw ($(mark_for_AD) + (0pt,-4pt)$) -- ($(mark_for_AD) + (0pt,4pt)$);
\node[blue] at ($(mark_for_AD) + (0pt,-8pt)$){$l_1$};
\coordinate (mark_for_BC) at ($(B)!0.5!(C)$);
\draw ($(mark_for_BC) + (0pt,-4pt)$) -- ($(mark_for_BC) + (0pt,4pt)$);
\node[red] at ($(mark_for_BC) + (0pt,8pt)$){$l_3$};

\coordinate (mark_for_AB) at ($(A)!0.5!(B)$);
\coordinate (mark_for_AB_up) at ($(mark_for_AB)! 1pt! (B)$);
\coordinate (mark_for_AB_down) at ($(mark_for_AB)! 1pt! (A)$);
\draw ($(mark_for_AB_up)!4pt!-90:(A)$) -- ($(mark_for_AB_up)!4pt!-90:(B)$);
\draw ($(mark_for_AB_down)!4pt!-90:(A)$) -- ($(mark_for_AB_down)!4pt!-90:(B)$);
\node[blue] at ($(mark_for_AB) + (-8pt,0pt)$){$l_2$};

\coordinate (mark_for_CD) at ($(C)!0.5!(D)$);
\coordinate (mark_for_CD_up) at ($(mark_for_CD)! 1pt! (C)$);
\coordinate (mark_for_CD_down) at ($(mark_for_CD)! 1pt! (D)$);
\draw ($(mark_for_CD_up)!4pt!-90:(D)$) -- ($(mark_for_CD_up)!4pt!-90:(C)$);
\draw ($(mark_for_CD_down)!4pt!-90:(D)$) -- ($(mark_for_CD_down)!4pt!-90:(C)$);
\node[red] at ($(mark_for_CD) + (8pt,0pt)$){$l_4$};

\end{tikzpicture}
\end{center}


Para resolver este ejercicio analizamos los datos que tenemos y cuales queremos obtener.\\
\textcolor{blue}{\underline{Datos:}} Recta $l_1$, recta $l_2$ y el punto medio de las diagonales.\\
\textcolor{red}{\underline{Objetivo:}} Encontrar las rectas $l_3$ y $l_4$.
        \vspace{5mm}

Como sabemos en un paralelógramo las caras opuestas tienen las mismas pendientes, en este caso $p_1=p_3$ y $p_2=p_3$.\\
Como primer paso obtendremos las pendientes de las rectas $l_1$ y $l_2$

\begin{minipage}{.4\textwidth}
\begin{align*}
 \dis  x+2y+4  &= 0\\
 \dis  y &= -\frac{x}{2}-2\\
 \dis  p_1 &= -\frac{1}{2}\\
\end{align*}

\end{minipage}% <---------------- Note the use of "%"
\begin{minipage}{.4\textwidth}

\begin{align*}
 \dis  5x+3y-1  &= 0\\
 \dis  y  &= -\frac{5x}{3}+\frac{1}{3}\\
 \dis  p_2 &= -\frac{5}{3}\\
\end{align*}

\end{minipage}

Para formar las rectas que estamos buscando demás de su pendiente necesitamos un punto que pertenezca a estas rectas. Para obtenerlo, primero buscamos el punto A intersectando las rectas $l_1$ y $l_2$. Luego obtendremos la distancia entre $A$
 y $E$, para luego utilizando la formula de distancia para encontrar el punto $C$.\\
 \vspace{5mm}
 
 Buscando $A$:
 
\begin{gather*}
x+2y+4 =  0\\
5x+3y-1 =  0
\end{gather*}
\begin{gather*}
5x+10y+20=  0 \\ 
5x+3y-1 =  0
\end{gather*}
\begin{gather*}
7y+21=  0 \\
y=-3
\end{gather*}
\begin{gather*}
x+2(-3)+4 =  0\\
x=2
\end{gather*}
Por lo tanto $A=(2,-3)$

 \vspace{5mm}
 
 
 Obtenemos la recta $\overline{AE}$:
 \begin{gather*}
(y-0)=\frac{0+3}{3-2}(x-3)\\
y=3x-9
\end{gather*}
 
 
  \vspace{5mm}
 
 Calculamos la distancia entre el punto $A=(2,-3)$ y $E=(3,0)$
\begin{align*}
D(A,E) &= \sqrt{(3-2)^2+(0+3)^2}\\
&= \sqrt{1+9} \\
&= \sqrt{10}
\end{align*}

Para encontrar el punto $C$ utilizamos la fórmula de distancia entre el punto $E$ y el punto $C$, cuya distancia es la misma que la del segmento $\overline{AE}$.

\begin{align*}
D(E,C) &= \sqrt{(3-x)^2+(0+y)^2}\\
\sqrt{10} &= \sqrt{(3-x)^2+(0+y)^2}\\
10 &= 9 - 6x + x^2 + y^2
\end{align*}
Ya que el punto que buscamos se encuentra en la misma recta que $\overline{AE}$, procedemos a reemplazar $y=3x-9$ para trabajar solo en una variable.

\begin{align*}
10 &= 9 - 6x + x^2 + (3x-9)^2\\
10 &= 9 - 6x + x^2 + 9x^2 - 54x + 81\\
0  &= 10x^2  - 60x + 80\\
0  &= x^2  - 6x + 8\\
0  &= (x-4)(x-2)\\
\end{align*}
Por lo tanto $x= 4$ y $x = 2$, pero como ya conocemos A que es $(2,-3)$ solo utilizaremos $x=4$ en la recta $y=3x-9$.
 \begin{gather*}
y=3x-9\\
y= 3(4)-9\\
y= 3
\end{gather*}
Entonces $C=(4,3)$

\vspace{5mm}

Finalmente buscaremos las rectas $l_3$ y $l_4$ utilizando $p_1=p_3=-\frac{1}{2}$, $p_2=p_4=-\frac{5}{3}$ y el punto $C=(4,3)$.\\

Recta $l_3$ usando $p_3=-\frac{1}{2}$ y el punto $C=(4,3)$:

 \begin{gather*}
l_3: y-3 = -\frac{1}{2} (x-4)\\
l_3: y= -\frac{x}{2} +5\\
\end{gather*}

Recta $l_4$ usando $p_3=-\frac{5}{3}$ y el punto $C=(4,3)$:

 \begin{gather*}
l_4: y-3 = -\frac{5}{3} (x-4)\\
l_4: y= -\frac{5x}{3} +\frac{20}{3}+3\\
l_4: y= -\frac{5x}{3} +\frac{29}{3}
\end{gather*}

\newpage

\item \textbf{Ejercicio:} Considere la ecuación $x^2-8x+y^2-6y+20=0$ y el punto $P = (8,1)$ \\
\begin{enumerate}
\item Determine el centro y el radio de la circunferencia.
\item Encuentre la ecuación principal de la recta que pasa por el centro de la circunferencia y el punto $P$
\item Halle los puntos de intersección entre la recta obtenida y la circunferencia dada. 
\end{enumerate}
\textbf{Solución:}

\begin{enumerate}
\item Encontramos el centro y radio de la circunferencia.
\begin{gather*}
x^2-8x+y^2-6y+20=0\\
x^2-8x+16-16+y^2-6y+9-9+20=0\\
x^2-8x+16+y^2-6y+9=5\\
(x-4)^2+(y-3)^2=5
\end{gather*}
El centro de la circunferencia esta en el punto $(4,3)$ y el radio es $\sqrt{5}$.

\item Recta que pasa por $P=(8,1)$ y el centro de la circunferencia $(4,3)$
\begin{gather*}
(y-1)= \frac{(1-3)}{(8-4)}(x-8)\\
y= -\frac{x}{2}+5
\end{gather*}

\begin{figure}[ht]
    \centering
    \includegraphics[width=15cm]{ayg1.png}
    \caption{Gráfica de la circunferencia y la recta obtenida}
    \label{ej3}
\end{figure}


\item Puntos de intersección entre la recta $y= -\frac{x}{2}+5$ y la circunferencia $x^2-8x+y^2-6y+20=0$

Para hacer un cálculo más simple utilizamos la función de la circunferencia factorizada $(x-4)^2+(y-3)^2=5$.

\begin{gather*}
(x-4)^2+(y-3)^2=5\\
(x-4)^2+((-\frac{x}{2}+5)-3)^2=5\\
(x-4)^2+(-\frac{x}{2}+2)^2=5\\
(x-4)^2+(\frac{4-x}{2})^2=5\\
(x-4)^2+\frac{(4-x)^2}{4}=5\\
(x-4)^2+\frac{(x-4)^2}{4}=5\\
\frac{5(x-4)^2}{4}=5\\
(x-4)^2=4\\
\abs{x-4}=2
\end{gather*}
Entonces $x=6$ y $x=2$.\\

Luego utilizando la recta $y= -\frac{x}{2}+5$ encontramos los puntos.
\begin{align*}
P_1 = (6,2)\\ 
P_2=(2,4)
\end{align*}
\end{enumerate}

\newpage

\item \textbf{Ejercicio:} Hallar la ecuación de la circunferencia que es tangente a la recta $l_1:x+y=3$ en el punto $(2,1)$ y cuyo centro se encuentra en la recta $3x-2y-6=0$\\

\textbf{Solución:}

Primero buscamos la recta perpendicular a $x+y=3$ que también se puede escribir $y=-x+3$.

\begin{gather*}
m_1 = -1\\
m_1 * m_2= -1\\
m_2 = 1
\end{gather*}
Hacemos la recta perpendicular con la pendiente $m_2= 1$ y el punto $(2,1)$.

\begin{gather*}
y-1= (1)(x-2)\\
y=x-1
\end{gather*}

Luego, el centro de la circunferencia lo encontramos en la interceccion de las rectas $3x-2y-6=0$ y $y=x-1$
\begin{gather*}
3x-2y-6=0\\
y=x-1
\end{gather*}

\begin{gather*}
3x-2(x-1)-6=0\\
3x-2x+2-6=0\\
x=4
\end{gather*}

\begin{gather*}
y=x-1\\
y=4-1\\
y=3
\end{gather*}

El centro de la circunferencia se encuentra en el punto $(4,3)$.\\

Para obtener el radio calculamos la distancia entre el punto $(4,3)$ y $(2,1)$

\begin{gather*}
D= \sqrt{(3-1)^2+(4-2)^2}\\
D= \sqrt{(2)^2+(2)^2}\\
D= \sqrt{8} = 2\sqrt{2}\\
\end{gather*}

Por lo tanto la ecuación de la circunferencia es:

\begin{gather*}
(x-4)^2+(y-3)^2=8
\end{gather*}


 \end{enumerate}


\end{document}
